\section{PHASE II: ZOOMING IN ON CELLS}
\label{sec:phase2}

Phase I provides each player with valid but wide confidence brackets on each cell's potential. This initial map is crucial for pruning the search space, but it is not sharp enough for final decision-making, primarily due to the center-versus-maximum bias. 

The purpose of Phase II is to resolve this ambiguity. Here, players ``zoom in" on the most promising regions identified in Phase I, conducting a localized exploration - which we call a ``local peek" - to refine their estimates and construct confidence bounds on the true cell maxima, $\mu^{\ast}(C)$, down to a pre-specified target accuracy, $\varepsilon > 0$.

The phase begins with each player independently forming a smaller ``active set" of candidate cells, $\Sact^{(0)}_j$. This is a safe elimination step: using their Phase I brackets, players discard any cell whose most optimistic outcome (its upper bound) cannot compete with the most pessimistic outcome (the lower bound) of the top-$N$ candidate cells. Players may form different active sets due to the randomness in their Phase I observations; our analysis handles this by considering the maximum active set size, $M_{act}:=\max_j |\Sact^{(0)}_j|$.

Within each of their active cells, players then conduct the local peek. For a fixed duration $T_1$, they sample from a fine $\eta$-net of probe points laid out inside each cell. An $\eta$-net is a grid of points so dense that any point in the cell is within a distance $\eta$ of some grid point. This strategy allows us to approximate the supremum of the Lipschitz function over the continuous cell by taking the maximum over a finite set of points. A probe at one of these points is successful only if no other player samples any point within the same cell in that round. While still subject to collisions, this exploration is highly focused on the regions that matter most.

\subsection{Collision-Tolerant Probe Sampling}
A key technical challenge is to ensure that, despite collisions and decentralization, every player gathers enough information at \emph{every} probe point. Our analysis shows that a carefully chosen phase duration $T_1$ is sufficient to guarantee this with high probability.

\begin{lemma}[Phase-II Probe Coverage]\label{lem:coverage}
Let $q_{M_{act},\eta}$ be the per-round success probability for a (player, cell, probe) triple, and let $\NumProbes$ be the total number of probe points across all players' active sets. Choose integers
\[
b \ge 4\log(2\NumProbes/\delta_{II}) \qquad\text{and}\qquad T_1 \ge 2b/q_{M_{act},\eta}.
\]
Then, with probability at least $1-\delta_{II}/2$, every triple attains at least $b$ non-collision samples.
\end{lemma}

With a guaranteed budget of at least $b$ successful samples per probe point, standard concentration inequalities ensure that the empirical mean at each point is a highly accurate estimate of its true mean. This accuracy at the probe points translates directly into a tight bound on the cell's maximum value.

\begin{proposition}[Refined Maxima Brackets]
\label{prop:refined-bracket}
At the end of Phase II, by choosing parameters $\eta$ and $b$ appropriately for a target accuracy $\varepsilon$, each player $j$ constructs a new, refined bracket $[\LCB^{(1)}_j(C), \UCB^{(1)}_j(C)]$ for each of her active cells. With high probability, this bracket contains the true maximum $\mu^\ast(C)$ and has a width of at most $\varepsilon$.
\end{proposition}

\subsection{Selecting the Top-$N$ Cells}
With these tight and reliable brackets on the true cell maxima, players are equipped to make their final selection. We adopt a deterministic rule to select $N$ cells: each player selects the $N$ cells corresponding to their largest lower confidence bounds, $\LCB^{(1)}_j(C)$, breaking any ties with a fixed, public ordering of the cells (e.g. say lexicographic). This guarantees every player outputs a set, $S^{(j)}_\varepsilon$, of size exactly $N$.

\begin{theorem}[Gap-Free $\varepsilon$-Optimality]\label{thm:gapfree-safety}
With high probability, the set $S^{(j)}_\varepsilon$ chosen by any player $j$ is $\varepsilon$-optimal: every cell $C \in S^{(j)}_\varepsilon$ satisfies $\mu^\ast(C) \ge \mu^\ast_{(N)} - \varepsilon$.
\end{theorem}

This powerful guarantee holds for any reward function, irrespective of the gaps between cell values. It ensures our procedure is robust even when the decision is difficult. Remarkably, this simple, decentralized rule leads to a powerful emergent behavior: under a mild separation condition on the reward function, it forces all players to agree on the exact same set of top-$N$ cells, achieving consensus without any communication.

\begin{definition}[$\varepsilon$-Uniqueness at the Top-$N$]
\label{def:eps-uniq}
A mean function $\mu$ is $\varepsilon$-unique at the top-$N$ if there is a unique set $S^\dagger \subset \Pcal$ of size $N$ such that $\min_{C\in S^\dagger} \mu^\ast(C) \ge \max_{C\notin S^\dagger} \mu^\ast(C) + 2\varepsilon$.
\end{definition}

\begin{lemma}[Consensus under $\varepsilon$-Uniqueness]\label{lem:consensus}
If the instance is $\varepsilon$-unique and the bracket width is at most $\varepsilon$, then on our high-probability event, all players select the exact same set of cells: $S^{(j)}_\varepsilon = S^\dagger$ for all $j \in [N]$.
\end{lemma}

This decentralized agreement is a cornerstone of our protocol's success since it allows the final seating phase to happen.

\begin{example}[Center-vs-maximum pathology in 1D]\label{ex:center-failure}
Let $d=1$, $h=\tfrac12$, so $\Pcal=\{C_1=[0,\tfrac12],\,C_2=[\tfrac12,1]\}$.
Define a Lipschitz mean $\mu(x)=x+\alpha \phi(x)$ with a narrow bump $\phi$ of height $1$ supported in $[\tfrac12-\delta,\tfrac12]$, with $\delta\ll h$ and $\alpha>0$ small so that $L$ is finite.
Then $\mu(x_{C_1})<\mu(x_{C_2})$ (center ranking favors $C_2$) but $\mu^\ast(C_1)>\mu^\ast(C_2)$ (the bump near the boundary makes $C_1$ optimal).
Phase~I center estimates therefore mis-rank the cells, while Phase~II's local peek in $C_1$ identifies the bump and restores the correct top-$N$ set.
\end{example}

\section{PHASE II\texorpdfstring{$\tfrac{1}{2}$}{1/2}: MUSICAL CHAIRS}
\label{sec:phase2half}

With a common set of $N$ high-quality target cells in hand, the players must perform the final coordination step: assigning themselves to these cells so that each is occupied by exactly one player. In our analysis, this common target set arises either from a consensus assumption, or from the public dither mechanism formalized in Appendix~G for the gap-free extension. This assignment is achieved via the Musical Chairs algorithm \citep{rosenski2016multi}. A key strength of our analysis is that we model the practical, challenging version of this protocol where players have no side-channel telling them which cells are already occupied; they must discover free cells through trial and error.

\subsection{The Seating Protocol}
The seating phase proceeds in rounds. Initially, all $N$ players are ``unseated." In each round, every unseated player samples a cell uniformly at random from the entire $N$-cell target set. A collision occurs at a cell if it is chosen by more than one unseated player, or if an unseated player chooses a cell that is already occupied by a seated player. If, however, a single unseated player chooses a currently unoccupied cell, that player becomes ``seated" at that cell. They cease to participate in the sampling process and will occupy that cell for the remainder of the horizon. The process terminates when all $N$ players are seated.

\subsection{Expected Seating Time and Regret}
Let $U_t$ be the number of unseated players at the start of round $t$. We define the expected number of players that become seated per round as "drift". When $u$ players remain unseated, this drift is given by $\Delta(u) := \EE[U_t - U_{t+1} \mid U_t=u] = \frac{u^2}{N}(1-\frac{1}{N})^{u-1}$. The quadratic dependence on $u$ means that progress is rapid when many players are searching for spots. This positive drift allows us to bound the total expected time for the dance to conclude.

\begin{theorem}[Expected Seating Time]
\label{thm:mc-time}
The expected time, $T_{MC}$, for all $N$ players to become seated is linear in the number of players: $\EE[T_{MC}] = O(N)$.
\end{theorem}

This result demonstrates that this simple, decentralized procedure is highly efficient and scales gracefully, far better than a naive $O(N \log N)$ coupon-collector analysis might suggest. The total regret incurred during this phase is therefore a fixed cost, dependent on $N$ but crucially, independent of the total time horizon $T$. This confirms that the entire coordination and seating process can be completed for a small, one-time price (not dependent on the time horizon).

\begin{corollary}[Seating-Phase Regret]
\label{cor:mc-regret}
The expected cumulative regret from Musical Chairs is bounded by a horizon-independent constant: $\EE[R_{MC}] = O(N^2)$.
\end{corollary}
